\documentclass[12pt,reqno]{amsart}
\usepackage{amsmath,amssymb,amsthm}
\usepackage{mathptmx}
\usepackage[scaled=0.90]{helvet}
\usepackage{courier}
\usepackage[T1]{fontenc}
\usepackage{enumitem}
\usepackage[hidelinks]{hyperref}
\pagestyle{plain}

\setlength{\oddsidemargin}{0cm}
\setlength{\evensidemargin}{0cm}
\setlength{\topmargin}{0in}
\setlength{\headheight}{0pt}
\setlength{\headsep}{0pt}
\setlength{\footskip}{.3in}
\setlength{\textheight}{9.2in}
\setlength{\textwidth}{6.5in}
\setlength{\parindent}{0pt}
\setlist[enumerate,1]{leftmargin=*,itemsep=1.1em}
\setlist[enumerate,2]{label=(\alph*),leftmargin=1.7em,itemsep=.45em,topsep=.4em}

\newcommand{\C}{\mathbb{C}}
\newcommand{\R}{\mathbb{R}}

\title{Problem Set 4 \\ MATH 411: Honors Complex Variables, Fall 2026}
\date{last compiled: \today}

% Coverage: Laurent series, isolated singularities, residues, and the
% residue theorem (Lectures 21--24), ahead of Midterm 2 on Fri, Nov. 6.
% Suggested completion date: Wednesday, November 4, 2026.
% Sources (paths relative to the course workspace):
% 1: moved from Problem Set 3, problem 7
%    (originally Old Homework/Homework/Homework5_Fall2021.tex, 3--4).
% 2--9: Old Homework/Homework/Homework5_Fall2021.tex, 2 and 5--12
%    (old 9 and 10 are combined as 6(a)--(b)).
% Homework5_Fall2021 problem 1 is omitted: it is already Problem Set 3,
% problem 6.
% Corrections to the 2021 version:
% - 3, last part: the 2021 statement printed a polynomial (residue 0); the
%   intended function, matching the 2021 answer 9/16, is
%   (z^3-1)(z+2)/(z^4-1)^2.
% - 4(b): the commented 2021 solution claimed the answer is 0. In fact
%   1 - sin z has double zeros at pi/2 + 2k pi, each with residue 2 for
%   the integrand; three lie inside |z| = 8, so the integral is 12 pi i.

\begin{document}
\maketitle

\begin{center}
Suggested completion date: Wednesday, November 4, 2026.
\end{center}

These are practice problems and are not to be turned in. Justify your
answers and all limiting arguments. Unless stated otherwise, closed
curves are oriented counterclockwise.

\begin{enumerate}

\item \textbf{Laurent expansions.}
\begin{enumerate}
\item Find a series in integer powers of $z$ representing
\[
 f(z)=\frac{1}{(z-1)(z-2)}
\]
in each of the regions $|z|<1$, $1<|z|<2$, and $|z|>2$.
Justify convergence in each region.
\item Find the first four nonzero terms of the Laurent expansion of
\[
 \frac{e^z}{z(1+z^2)}
\]
valid for $0<|z|<1$.
\end{enumerate}

\item \textbf{Logarithmic derivatives.}
\begin{enumerate}
\item Show that the association $f\mapsto f'/f$ (for $f$ holomorphic
and nonvanishing) sends products to sums.
\item If $P(z)=(z-a_1)\cdots(z-a_n)$, where $a_1,\dots,a_n\in\C$ are
the roots of $P$ (listed with multiplicity), compute $P'/P$.
\item Let $\gamma$ be a closed path on which $P$ has no roots. Show
that
\[
 \frac{1}{2\pi i}\int_\gamma\frac{P'(z)}{P(z)}\,dz
 = W(\gamma,a_1)+\cdots+W(\gamma,a_n).
\]
\end{enumerate}

\item \textbf{Residues at the origin.}
Find the residue at $0$ of each of the following.
\begin{enumerate}
\item $\dfrac{2z+1}{z(z^3-5)}$
\item $\dfrac{e^z}{z^4}$
\item $\dfrac{e^z}{\sin z}$
\end{enumerate}
Also find the residue at $z=1$ of
\[
 \frac{(z^3-1)(z+2)}{(z^4-1)^2}.
\]

\item \textbf{Contour integrals by residues.}
\begin{enumerate}
\item Evaluate
\[
 \int_C\frac{dz}{z^2-3z+5},
\]
where $C$ is the boundary of the rectangle with vertices
$0,\ 10,\ 10+4i,\ 4i$, oriented \emph{clockwise}.
\item Evaluate
\[
 \int_C\frac{1+z}{1-\sin z}\,dz,
\]
where $C$ is the circle $|z|=8$. \emph{Hint:} First find all complex
solutions of $\sin z=1$ and determine the order of each zero of
$1-\sin z$.
\end{enumerate}

\item \textbf{A trigonometric integral.}
Show that
\[
 \int_0^{2\pi}\frac{dt}{(5-3\sin t)^2}=\frac{5\pi}{32}.
\]

\item \textbf{Rational integrals over the real line.}
\begin{enumerate}
\item Show that
\[
 \int_0^\infty\frac{x^2}{x^6+1}\,dx=\frac{\pi}{6}.
\]
\item For $a,b>0$, show that
\[
 \int_0^\infty\frac{dx}{(x^2+a^2)(x^2+b^2)}=\frac{\pi}{2ab(a+b)}.
\]
\end{enumerate}

\item \textbf{A Fourier-type integral.}
For $a,b>0$ with $a\neq b$, show that
\[
 \int_{-\infty}^{\infty}\frac{\cos x}{(x^2+a^2)(x^2+b^2)}\,dx
 =\frac{\pi}{a^2-b^2}\left(\frac{e^{-b}}{b}-\frac{e^{-a}}{a}\right).
\]

\item \textbf{An integral with a logarithm.}
Show that
\[
 \int_0^\infty\frac{(\log x)^2}{1+x^2}\,dx=\frac{\pi^3}{8}.
\]
\emph{Hint:} Use a branch of $\log z$ on $\C$ minus the negative
imaginary axis, and integrate $(\log z)^2/(1+z^2)$ around the boundary
of the region $\{\delta\leq|z|\leq R,\ \operatorname{Im}z\geq0\}$
(Lang, \S VI.2, Exercise 14(a)). You may use
$\int_0^\infty dx/(1+x^2)=\pi/2$.

\item \textbf{The Gaussian integral via residues.}
Prove that
\[
 \int_{-\infty}^{\infty}e^{-x^2}\,dx=\sqrt{\pi}
\]
by integrating
\[
 f(z)=\frac{e^{-z^2}}{1+e^{-2az}},\qquad a=e^{\pi i/4}\sqrt{\pi},
\]
around the parallelogram with vertices $-R,\ R,\ R+a,\ -R+a$ and then
letting $R\to\infty$. \emph{Hint:} Check that $f(z)-f(z+a)=e^{-z^2}$.

\end{enumerate}
\end{document}
